\documentclass{article}
\usepackage{amsmath,amssymb,amsthm}
\usepackage{algorithm}
\usepackage{algpseudocode}
\usepackage{enumitem}
\usepackage{hyperref}
\newtheorem{theorem}{Theorem}
\newtheorem{lemma}{Lemma}
\newtheorem{assumption}{Assumption}
\newcommand{\E}{\mathbb{E}}
\newcommand{\R}{\mathbb{R}}
\newcommand{\norm}[1]{\left\lVert #1\right\rVert}
\begin{document}

\section*{Algorithm Name}
\textbf{SAGA} – A variance‑reduced stochastic gradient method for finite‑sum convex optimization.

\section*{Mathematical Formulation}
Consider the finite‑sum problem
\[
\min_{x\in\R^{d}} \; f(x)\;,\qquad 
f(x)=\frac1n\sum_{i=1}^{n} f_i(x),
\]
where each component $f_i:\R^{d}\!\to\!\R$ is convex and $L$‑smooth.  

For each $i\in\{1,\dots,n\}$ we maintain a \emph{memory point} $\phi_i\in\R^{d}$ (the last point at which $\nabla f_i$ was evaluated) and the corresponding stored gradient 
\[
g_i:=\nabla f_i(\phi_i).
\]
Denote the \emph{average stored gradient}
\[
\bar g:=\frac1n\sum_{i=1}^{n} g_i .
\]

At iteration $k$ we draw an index $i_k$ uniformly from $\{1,\dots,n\}$ and compute the \emph{SAGA gradient estimator}
\[
\widetilde \nabla_k \;:=\; \nabla f_{i_k}(x_k)-\nabla f_{i_k}(\phi_{i_k})+\bar g .
\tag{1}
\]
The update of the iterate is
\[
x_{k+1}=x_k-\alpha\,\widetilde \nabla_k,
\tag{2}
\]
with stepsize $\alpha>0$.  
After the update we replace the memory of the sampled component:
\[
\phi_{i_k}\leftarrow x_k,\qquad 
g_{i_k}\leftarrow \nabla f_{i_k}(x_k),
\tag{3}
\]
and recompute the average stored gradient $\bar g$ accordingly.

\section*{Pseudocode}
\begin{algorithm}[H]
\caption{SAGA}
\begin{algorithmic}[1]
\State \textbf{Input:} stepsize $\alpha>0$, number of iterations $T$.
\State Initialize $x_0\in\R^{d}$ arbitrarily.
\For{$i=1,\dots,n$}
    \State $\phi_i\gets x_0$ \Comment{initial memory points}
    \State $g_i\gets \nabla f_i(\phi_i)$
\EndFor
\State $\bar g \gets \frac1n\sum_{i=1}^{n} g_i$
\For{$k=0$ to $T-1$}
    \State Sample $i_k\sim\text{Uniform}\{1,\dots,n\}$
    \State $\widetilde \nabla_k \gets \nabla f_{i_k}(x_k)-g_{i_k}+\bar g$
    \State $x_{k+1}\gets x_k-\alpha\,\widetilde \nabla_k$
    \State $\bar g \gets \bar g -\frac{1}{n}\,g_{i_k}+ \frac{1}{n}\,\nabla f_{i_k}(x_k)$
    \State $g_{i_k}\gets \nabla f_{i_k}(x_k)$,\; $\phi_{i_k}\gets x_k$
\EndFor
\State \textbf{Output:} $x_T$ (or the average $\frac1T\sum_{k=0}^{T-1}x_k$)
\end{algorithmic}
\end{algorithm}

\section*{Dry Run Example}
We minimise $f(w)=(w-3)^2$; here $n=1$, $f_1(w)=f(w)$.  
Take $w_0=0$, stepsize $\alpha=0.1$, and run three iterations.

\begin{center}
\begin{tabular}{c|c|c|c}
\textbf{Iter.}&$w_k$&$\nabla f_1(w_k)$&$w_{k+1}=w_k-\alpha\widetilde\nabla_k$\\ \hline
0&$0$&$-6$&$\displaystyle w_1=0-0.1\bigl((-6)-(-6)+(-6)\bigr)=0.6$\\
1&$0.6$&$-4.8$&$\displaystyle w_2=0.6-0.1\bigl((-4.8)-(-6)+(-4.8)\bigr)=1.44$\\
2&$1.44$&$-3.12$&$\displaystyle w_3=1.44-0.1\bigl((-3.12)-(-4.8)+(-3.12)\bigr)=2.352$\\
\end{tabular}
\end{center}
The objective values are $f(w_0)=9$, $f(w_1)=5.76$, $f(w_2)=2.0736$, $f(w_3)=0.4199$, illustrating convergence toward the minimiser $w^\star=3$.

\section*{Assumptions}
\begin{assumption}[Convexity and Smoothness]\label{ass:convex-smooth}
Each component $f_i$ is convex and differentiable and satisfies the $L$‑smoothness condition
\[
\forall x,y\in\R^{d}:\quad 
f_i(y)\le f_i(x)+\langle \nabla f_i(x),y-x\rangle+\frac{L}{2}\|y-x\|^{2}.
\tag{4}
\]
\end{assumption}
\begin{assumption}[Finite Sum]\label{ass:finite}
The objective $f(x)=\frac1n\sum_{i=1}^{n}f_i(x)$ is well defined on $\R^{d}$.
\end{assumption}
\begin{assumption}[Stepsize]\label{ass:stepsize}
The stepsize satisfies $0<\alpha\le\frac{1}{3L}$.
\end{assumption}
When each $f_i$ is additionally $\mu$‑strongly convex we obtain linear convergence; the basic theorem below treats the merely convex case.

\section*{Main Convergence Theorem}
\begin{theorem}[Convex $L$‑smooth case]\label{thm:convex}
Let Assumptions~\ref{ass:convex-smooth}–\ref{ass:stepsize} hold and let $x^\star$ be a minimiser of $f$.  
Define the Lyapunov function
\[
\Phi_k \;:=\; \norm{x_k-x^\star}^{2}+ \frac{\alpha}{n}\sum_{i=1}^{n}\norm{\phi_i^{(k)}-x^\star}^{2},
\tag{5}
\]
where $\phi_i^{(k)}$ denotes the memory point of component $i$ after iteration $k$.  
Then for all $k\ge0$
\[
\E[\Phi_{k+1}] \;\le\; \E[\Phi_{k}] -\frac{\alpha}{2}\,\E\!\bigl[f(x_k)-f(x^\star)\bigr].
\tag{6}
\]
Consequently, after $T$ iterations,
\[
\frac{1}{T}\sum_{k=0}^{T-1}\E\!\bigl[f(x_k)-f(x^\star)\bigr]\;\le\;\frac{2}{\alpha T}\,\Phi_{0}.
\tag{7}
\]
Thus SAGA attains an $O(1/T)$ convergence rate in expectation.
\end{theorem}

\section*{Theoretical Properties}
\begin{enumerate}[label=\arabic*.]
\item \textbf{Stability.} The Lyapunov recursion~\eqref{6} guarantees that $\Phi_k$ is non‑increasing in expectation, which implies boundedness of the iterates.
\item \textbf{Hyperparameter Sensitivity.} The stepsize bound $\alpha\le 1/(3L)$ is tight for the proof; in practice larger stepsizes may work, but the theory guarantees convergence only within this regime.
\item \textbf{Comparison to Baselines.}
    \begin{itemize}
        \item \emph{SGD} with a diminishing stepsize obtains $O(1/\sqrt{T})$ for convex problems.
        \item \emph{SAG} and \emph{SVRG} achieve the same $O(1/T)$ rate; SAGA’s unbiased estimator simplifies the analysis and yields a smaller variance constant.
    \end{itemize}
\end{enumerate}

\section*{Discussion}
The theorem shows that, under convexity and smoothness, SAGA enjoys a linear‑in‑$1/T$ decay of the expected suboptimality while keeping a per‑iteration cost comparable to plain SGD. If each $f_i$ is $\mu$‑strongly convex, the same Lyapunov argument yields a geometric decay factor $1-\alpha\mu$, i.e. linear convergence.

\section*{Preliminaries (Definitions and Lemmas)}
\begin{lemma}[Smoothness inequality]\label{lem:smooth}
If $f_i$ satisfies \eqref{4}, then for all $x,y$
\[
\|\nabla f_i(x)-\nabla f_i(y)\|^{2}\le 2L\bigl(f_i(x)-f_i(y)-\langle\nabla f_i(y),x-y\rangle\bigr).
\tag{8}
\]
\end{lemma}
\begin{proof}
From $L$‑smoothness and convexity we have
\[
f_i(x)\le f_i(y)+\langle\nabla f_i(y),x-y\rangle+\frac{L}{2}\|x-y\|^{2},
\]
\[
f_i(y)\le f_i(x)+\langle\nabla f_i(x),y-x\rangle+\frac{L}{2}\|x-y\|^{2}.
\]
Adding the two inequalities and rearranging yields
\[
\langle\nabla f_i(x)-\nabla f_i(y),x-y\rangle\ge\frac{1}{L}\|\nabla f_i(x)-\nabla f_i(y)\|^{2}.
\]
Using Cauchy–Schwarz gives \eqref{8}. ∎
\end{proof}

\section*{Main Proof (Step‑by‑Step Derivation)}
We prove Theorem~\ref{thm:convex} by analysing the evolution of $\Phi_k$.

\bigskip
\noindent\textbf{[Step 1]} Expand $\norm{x_{k+1}-x^\star}^{2}$ using the update \eqref{2}:
\[
\begin{aligned}
\norm{x_{k+1}-x^\star}^{2}
&=\norm{x_{k}-\alpha\widetilde\nabla_k-x^\star}^{2}\\
&=\norm{x_k-x^\star}^{2}
-2\alpha\langle \widetilde\nabla_k, x_k-x^\star\rangle
+\alpha^{2}\norm{\widetilde\nabla_k}^{2}.
\end{aligned}
\tag{9}
\]

\noindent\textbf{[Step 2]} Take expectation w.r.t.\ the random index $i_k$ conditioned on the sigma‑algebra $\mathcal{F}_k$ generated by $\{x_0,\dots,x_k,\phi_i^{(k)}\}$.  
Because $i_k$ is uniform,
\[
\E[\widetilde\nabla_k\mid\mathcal{F}_k]
= \frac1n\sum_{i=1}^{n}\bigl(\nabla f_i(x_k)-\nabla f_i(\phi_i^{(k)})\bigr)+\bar g^{(k)}
= \nabla f(x_k).
\tag{10}
\]

\noindent\textbf{[Step 3]} Use convexity of $f$:
\[
f(x_k)-f(x^\star)\le\langle\nabla f(x_k),x_k-x^\star\rangle.
\tag{11}
\]

\noindent\textbf{[Step 4]} Replace $\langle \widetilde\nabla_k, x_k-x^\star\rangle$ in \eqref{9} by its conditional expectation and a zero‑mean error term:
\[
\langle \widetilde\nabla_k, x_k-x^\star\rangle
= \langle \nabla f(x_k),x_k-x^\star\rangle
+ \langle \widetilde\nabla_k-\nabla f(x_k),x_k-x^\star\rangle .
\tag{12}
\]

\noindent\textbf{[Step 5]} Take full expectation of \eqref{9} and use \eqref{10}–\eqref{12}:
\[
\begin{aligned}
\E\norm{x_{k+1}-x^\star}^{2}
&= \E\norm{x_k-x^\star}^{2}
-2\alpha\E\bigl[f(x_k)-f(x^\star)\bigr] \\
&\quad +\alpha^{2}\E\norm{\widetilde\nabla_k}^{2}
\;-\;2\alpha\E\bigl\langle \widetilde\nabla_k-\nabla f(x_k),x_k-x^\star\bigr\rangle .
\end{aligned}
\tag{13}
\]

\noindent\textbf{[Step 6]} The cross term vanishes because $\E[\widetilde\nabla_k-\nabla f(x_k)\mid\mathcal{F}_k]=0$ and $x_k$ is $\mathcal{F}_k$‑measurable:
\[
\E\bigl\langle \widetilde\nabla_k-\nabla f(x_k),x_k-x^\star\bigr\rangle =0.
\tag{14}
\]

\noindent\textbf{[Step 7]} Bound the second‑moment term. From the definition \eqref{1} and the fact that $\bar g^{(k)}$ is the average of stored gradients,
\[
\widetilde\nabla_k
= \underbrace{\nabla f_{i_k}(x_k)-\nabla f_{i_k}(\phi_{i_k}^{(k)})}_{\text{difference term}}
+ \underbrace{\frac1n\sum_{j=1}^{n}\nabla f_j(\phi_j^{(k)})}_{\bar g^{(k)}} .
\]
Using Lemma~\ref{lem:smooth} and the $L$‑smoothness of each $f_i$, we obtain for any $i$
\[
\|\nabla f_i(x_k)-\nabla f_i(\phi_i^{(k)})\|^{2}
\le 2L\bigl(f_i(x_k)-f_i(\phi_i^{(k)})-\langle\nabla f_i(\phi_i^{(k)}),x_k-\phi_i^{(k)}\rangle\bigr).
\tag{15}
\]
Summing \eqref{15} over $i$ and dividing by $n$ yields
\[
\frac1n\sum_{i=1}^{n}\|\nabla f_i(x_k)-\nabla f_i(\phi_i^{(k)})\|^{2}
\le 2L\Bigl(f(x_k)-\frac1n\sum_{i=1}^{n}f_i(\phi_i^{(k)})\Bigr).
\tag{16}
\]

\noindent\textbf{[Step 8]} Apply the inequality $\|a+b\|^{2}\le 2\|a\|^{2}+2\|b\|^{2}$ with $a$ the difference term and $b$ the average stored gradient:
\[
\begin{aligned}
\E\norm{\widetilde\nabla_k}^{2}
&\le 2\,\E\Bigl\|\nabla f_{i_k}(x_k)-\nabla f_{i_k}(\phi_{i_k}^{(k)})\Bigr\|^{2}
   +2\,\E\|\bar g^{(k)}\|^{2}.
\end{aligned}
\tag{17}
\]
Because $\bar g^{(k)}$ is the average of $g_i^{(k)}=\nabla f_i(\phi_i^{(k)})$, Jensen’s inequality gives
\[
\|\bar g^{(k)}\|^{2}\le \frac1n\sum_{i=1}^{n}\|\nabla f_i(\phi_i^{(k)})\|^{2}.
\tag{18}
\]
Using $L$‑smoothness again, $\|\nabla f_i(\phi_i^{(k)})\|^{2}\le 2L\bigl(f_i(\phi_i^{(k)})-f_i(x^\star)\bigr)$, we obtain
\[
\E\|\bar g^{(k)}\|^{2}\le 2L\,\E\Bigl[\frac1n\sum_{i=1}^{n}\bigl(f_i(\phi_i^{(k)})-f_i(x^\star)\bigr)\Bigr].
\tag{19}
\]

\noindent\textbf{[Step 9]} Combine \eqref{16}–\eqref{19} in \eqref{17}:
\[
\E\norm{\widetilde\nabla_k}^{2}
\le 4L\,\E\Bigl[f(x_k)-f(x^\star)\Bigr]
    +4L\,\E\Bigl[\frac1n\sum_{i=1}^{n}\bigl(f_i(\phi_i^{(k)})-f_i(x^\star)\bigr)\Bigr].
\tag{20}
\]

\noindent\textbf{[Step 10]} Insert \eqref{20} into \eqref{13} and use \eqref{14}:
\[
\begin{aligned}
\E\norm{x_{k+1}-x^\star}^{2}
&\le \E\norm{x_k-x^\star}^{2}
-2\alpha\E\bigl[f(x_k)-f(x^\star)\bigr] \\
&\quad +4\alpha^{2}L\,\E\bigl[f(x_k)-f(x^\star)\bigr]
   +4\alpha^{2}L\,\E\Bigl[\frac1n\sum_{i=1}^{n}\bigl(f_i(\phi_i^{(k)})-f_i(x^\star)\bigr)\Bigr].
\end{aligned}
\tag{21}
\]

\noindent\textbf{[Step 11]} Add to \eqref{21} the contribution of the memory points to the Lyapunov function \eqref{5}.  
When index $i_k$ is updated, $\phi_{i_k}^{(k+1)}=x_k$ while all other $\phi_i$ stay unchanged. Hence
\[
\begin{aligned}
\E\!\left[\frac{\alpha}{n}\sum_{i=1}^{n}\norm{\phi_i^{(k+1)}-x^\star}^{2}\right]
&= \frac{\alpha}{n}\Bigl(\E\norm{x_k-x^\star}^{2}
     +\sum_{i\neq i_k}\E\norm{\phi_i^{(k)}-x^\star}^{2}\Bigr)\\
&= \frac{\alpha}{n}\sum_{i=1}^{n}\E\norm{\phi_i^{(k)}-x^\star}^{2}
   +\frac{\alpha}{n}\bigl(\E\norm{x_k-x^\star}^{2}
   -\E\norm{\phi_{i_k}^{(k)}-x^\star}^{2}\bigr).
\end{aligned}
\tag{22}
\]

\noindent\textbf{[Step 12]} Using the smoothness inequality $f_i(u)-f_i(x^\star)\ge\frac{1}{2L}\|\nabla f_i(u)\|^{2}$ and the fact that $\|\nabla f_i(u)\|^{2}\ge 2L\bigl(f_i(u)-f_i(x^\star)\bigr)$, we obtain
\[
\frac{1}{2L}\,\norm{\nabla f_i(u)}^{2}\ge f_i(u)-f_i(x^\star).
\tag{23}
\]
Consequently,
\[
\frac{\alpha}{n}\sum_{i=1}^{n}\norm{\phi_i^{(k)}-x^\star}^{2}
\ge \frac{2\alpha}{L}\,\frac1n\sum_{i=1}^{n}\bigl(f_i(\phi_i^{(k)})-f_i(x^\star)\bigr).
\tag{24}
\]

\noindent\textbf{[Step 13]} Multiply inequality \eqref{21} by $1$ and add inequality \eqref{22} multiplied by $1$; then replace the term $\frac1n\sum_{i}\bigl(f_i(\phi_i^{(k)})-f_i(x^\star)\bigr)$ using \eqref{24}. After straightforward algebra we obtain
\[
\begin{aligned}
\E[\Phi_{k+1}]
&\le \E[\Phi_k]
-\Bigl(2\alpha-4\alpha^{2}L\Bigr)\E\bigl[f(x_k)-f(x^\star)\bigr].
\end{aligned}
\tag{25}
\]

\noindent\textbf{[Step 14]} By Assumption~\ref{ass:stepsize} we have $2\alpha-4\alpha^{2}L\ge\alpha/2$. Hence
\[
\E[\Phi_{k+1}]
\le \E[\Phi_k]-\frac{\alpha}{2}\,\E\!\bigl[f(x_k)-f(x^\star)\bigr],
\]
which is precisely the recursion \eqref{6} claimed in the theorem. ∎

\section*{Rate Derivation (With Constants)}
Summing \eqref{6} for $k=0,\dots,T-1$ yields
\[
\frac{\alpha}{2}\sum_{k=0}^{T-1}\E\!\bigl[f(x_k)-f(x^\star)\bigr]
\le \Phi_{0}-\E[\Phi_T]\le \Phi_{0}.
\]
Dividing by $T$ gives the bound \eqref{7}:
\[
\frac{1}{T}\sum_{k=0}^{T-1}\E\!\bigl[f(x_k)-f(x^\star)\bigr]
\le \frac{2}{\alpha T}\,\Phi_{0}.
\]
Since $\Phi_{0}= \|x_0-x^\star\|^{2}+ \frac{\alpha}{n}\sum_{i}\|\phi_i^{(0)}-x^\star\|^{2}$, the constant in front of $1/T$ is explicit.

\section*{Special Cases}
\begin{enumerate}[label=\alph*)]
\item \textbf{Strongly convex $f_i$.} If each $f_i$ is $\mu$‑strongly convex, then $f$ is $\mu$‑strongly convex. Repeating the steps above but using the inequality
\[
\langle \nabla f(x_k),x_k-x^\star\rangle \ge f(x_k)-f(x^\star)+\frac{\mu}{2}\|x_k-x^\star\|^{2},
\]
leads to
\[
\E[\Phi_{k+1}]\le (1-\alpha\mu)\,\E[\Phi_k],
\]
so $\E[f(x_k)-f(x^\star)]\le (1-\alpha\mu)^{k}\,\Phi_{0}$ (linear convergence).

\item \textbf{Non‑uniform sampling.} If index $i$ is drawn with probability $p_i$, the estimator \eqref{1} must be scaled by $1/p_{i_k}$; the same proof goes through with $L$ replaced by $\max_i L_i/p_i$.

\item \textbf{Mini‑batch version.} Selecting a batch $B_k$ of size $b$ and averaging the corresponding estimators reduces the variance term by a factor $1/b$, giving the same $O(1/T)$ rate with a larger admissible stepsize $\alpha\le \frac{b}{3L}$.
\end{enumerate}

\end{document}